% !TeX spellcheck = it_IT
\documentclass[a4paper,10pt]{article}
\usepackage[utf8]{inputenc}

% ----  Useful packages % ---- 
\usepackage{amsmath}
\usepackage{graphicx}
\usepackage{amsfonts}
\usepackage{amsthm}
\usepackage{amssymb}
\usepackage{makecell}
\usepackage{array}
\usepackage{booktabs}
\usepackage{multirow}
% ----  Useful packages % ---- 

\usepackage{wrapfig}
\usepackage{caption}
\usepackage{subcaption}
\usepackage{hyperref}
\hypersetup{
    colorlinks,
    citecolor=black,
    filecolor=black,
    linkcolor=black,
    urlcolor=black
}

\graphicspath{ {./images/} }

% ---- Set page size and margins replace ------
\usepackage[letterpaper,top=2cm,bottom=2cm,left=3cm,right=3cm,marginparwidth=1.75cm]{geometry}
% ---- Set page size and margins replace ------

% ------- NOTA ------
\theoremstyle{remark}
\newtheorem{note}{Note}[subsubsection]
% ------- NOTA ------

% ------- OSSERVAZIONE ------
\theoremstyle{definition}
\newtheorem{observation}{Osservazione}[subsection]
% ------- OSSERVAZIONE ------

% ------- DEFINIZIONE ------
\theoremstyle{plain}
\newtheorem{definition}{Definizione}[subsection]
% ------- DEFINIZIONE ------

% ------- ESEMPIO ------
\theoremstyle{definition}
\newtheorem{example}{Esempio}[subsection]
% ------- ESEMPIO ------

% ------- DIMOSTRAZIONE ------
\theoremstyle{definition}
\newtheorem{demostration}{Dimostrazione}[subsection]
% ------- DIMOSTRAZIONE ------

% ------- TEOREMA ------
\theoremstyle{definition}
\newtheorem{theorem}{Teorema}[subsection]
% ------- TEOREMA ------

% ------- COROLLARIO ------
\theoremstyle{plain}
\newtheorem{corollaries}{Corollario}[theorem]
% ------- COROLLARIO ------

% ------- PROPOSIZIONE ------
\theoremstyle{plain}
\newtheorem{proposition}{Proposizione}[subsection]
% ------- PROPOSIZIONE ------

% ---- Footer and header ---- 
\usepackage{fancyhdr}
\pagestyle{fancy}
\fancyhf{}
\fancyhead[LE,RO]{A.A 2023-2024}
\fancyhead[RE,LO]{Laboratorio II}
\fancyfoot[RE,LO]{\rightmark}
\fancyfoot[LE,RO]{\thepage}

\renewcommand{\headrulewidth}{.5pt}
\renewcommand{\footrulewidth}{.5pt}
% ---- Footer and header ---- 

% ----  Language setting ---- 
\usepackage[italian, english]{babel}
% ----  Language setting ---- 

\usepackage{listings}
\usepackage{color}

\definecolor{dkgreen}{rgb}{0,0.6,0}
\definecolor{gray}{rgb}{0.5,0.5,0.5}
\definecolor{mauve}{rgb}{0.58,0,0.82}

\lstset{frame=tb,
  language=C,
  aboveskip=3mm,
  belowskip=3mm,
  showstringspaces=false,
  columns=flexible,
  basicstyle={\small\ttfamily},
  numbers=none,
  numberstyle=\tiny\color{gray},
  keywordstyle=\color{blue},
  commentstyle=\color{dkgreen},
  stringstyle=\color{mauve},
  breaklines=true,
  breakatwhitespace=true,
  tabsize=3
}

\begin{document}
\begin{titlepage} %crea l'enviroment
	\begin{figure}[t] %inserisce le figure
		\centering\includegraphics[width=0.98\textwidth]{marchio_unipi_pant541.png}
	\end{figure}
	\vspace{20mm}
	
	\begin{Large}
		\begin{center}
			\textbf{Dipartimento di Informatica\\ Corso di Laurea Triennale in Informatica\\}
			\vspace{20mm}
			{\LARGE{Corso a Libera Scelta - 6 CFU}}\\
			\vspace{10mm}
			{\huge{\bf Computer Graphics}}\\
		\end{center}
	\end{Large}
	
	
	\vspace{36mm}
	%minipage divide la pagina in due sezioni settabili
	\begin{minipage}[t]{0.47\textwidth}
		{\large{\bf Professore:}\\ \large{Prof. }}
	\end{minipage}
	\hfill
	\begin{minipage}[t]{0.47\textwidth}\raggedleft
		{\large{\bf Autore:}\\ \large{Filippo Ghirardini}}
	\end{minipage}
	
	\vspace{25mm}
	
	\hrulefill
	
	\vspace{5mm}
	
	\centering{\large{\bf Anno Accademico 2023/2024 }}
	
\end{titlepage}

\tableofcontents
\begin{center}
    \vspace{-20pt}
    \rule{11cm}{.1pt} 
\end{center}
\newpage
\section{Funzioni specifiche}
\subsection{perror}
\textbf{Come funziona \textit{perror}?}\\
È una funzione che stampa un messaggio di errore su \textit{stderr} seguito da una descrizione. Si basa sulla variabile globale \textit{errno}.
\begin{lstlisting}[language=C, mathescape]
	void perror(const char *msg) $\to$ msg:string
\end{lstlisting}
\subsection{stderr e stdout}
\textbf{Qual è l'utilizzo di \textit{stderr} e \textit{stdout}?}\\
\textit{stdout} lo usiamo per l'output, di default i dati vengono stampati sul terminare (e.g. \textit{printf}). Per ridirezionarlo usiamo
\begin{lstlisting}[language=BASH]
	command > file_output
\end{lstlisting}
\textit{stderr} è usato invece per l'output degli errori o messaggi diagnostici (e.g. \textit{perror}). Per ridirezionarlo usiamo
\begin{lstlisting}[language=BASH]
	command >&2 file_output
\end{lstlisting}
\subsection{fscanf}
\textbf{Cosa significano i modificatori \textit{\%NS} e \textit{\%ms} in \textit{fscanf}?}\\
\begin{itemize}
	\item \textbf{\%Ns} è un modificatore che permette di leggere $d$ massimo $n-1$ caratteri in input, riservando il carattere finale \textit{$\backslash 0$}.
	\begin{lstlisting}[language=C]
		fscanf(file, "%Ns", buffer);
	\end{lstlisting}
	Con \textit{buffer} già allocato.
	\item \textbf{\%ms} è un modificatore che permette di allocare dinamicamente la memoria necessaria per memorizzare la stringa.
	\begin{lstlisting}[language=C]
		fscanf(file, "%ms", &buffer);
	\end{lstlisting}
	Con \textit{buffer} un puntatore a \textit{char*}. La memoria è deallocata automaticamente.
\end{itemize}
\subsection{snprintf e asprintf}
\textbf{Allocazione di stringhe tramite \textit{snprintf} e \textit{asprintf}.}\\
Con \textit{snprintf} si costruiscono stringhe in modo sicuro scrivendole in un buffer preallocato mentre \textit{asprintf} alloca dinamicamente la memoria.
\begin{itemize}
	\item \begin{lstlisting}[language=C]
		int snprintf(char* str, size_t size, const char* format);
	\end{lstlisting}
	Questa funzione restituisce il numero di caratteri scritti e usa i seguenti parametri:
	\begin{itemize}
		\item \textit{str}: puntatore al buffer dove verrà scritta la stringa
		\item \textit{size}: dimensione massima del buffer incluso \textit{$\backslash 0$}
		\item \textit{format}: stringa di formato
	\end{itemize}
	\item \begin{lstlisting}[language=C, mathescape]
		int asprintf(char ** strp, const char * format, $\ldots$);
	\end{lstlisting}
	Restituisce il numero di caratteri scritti senza $\backslash 0$ e prende come parametro \textit{strp} ovvero il puntatore dove verrà allocata la memoria per la stringa.
\end{itemize}
\subsection{fscanf}
\textbf{Quali sono le limitazioni di \textit{fscanf}?}\\
\begin{itemize}
	\item È sensibile al formato, quindi se passo \textit{\%d} e poi leggo un intero mi restituisce errore
	\item Non c'è la gestione delle righe, ignora i caratteri \textit{$\backslash n$}
	\item Legge fino al prossimo spazio, può causare \textbf{troncamenti}
	\item Non gestisce facilmente le virgole e non verifica i limiti dei buffer
\end{itemize}
\subsection{getline}
\textbf{Spiega la funzione \textit{getline}.}\\
È una funzione per leggere intere righe di testo.
\begin{lstlisting}[language=C]
	ssize_t	 getline(char ** lineptr, size_t * n, FILE* stream);
\end{lstlisting}
Dove i parametri sono:
\begin{itemize}
	\item \textit{lineptr}: punta alla memoria dove verrà memorizzata la riga letta. Se è \textit{NULL} la funzione alloca dinamicamente. Se è troppo piccolo il buffer viene riallocato.
	\item \textit{n}: la dimensione del buffer, viene aggiornato automaticamente se viene ridimensionato
	\item \textit{stream}: puntatore al file da cui leggere
\end{itemize}
\subsection{fscanf vs getline}
\textbf{Qual è la differenza sostanziale tra \textit{fscanf} e \textit{getline}?}\\
\textit{fscanf} è ottima per leggere dati formattati direttamente ma richiede attenzione ai rischi di buffer overflow. \textit{getline} è più flessibile e sicuro per leggere righe intere ma necessita di ulteriori elaborazioni per estrarre i dati interni.
\subsection{strtok}
\textbf{Parsing di stringhe con \textit{strtok}.}\\
La funzione \textit{strtok} è usata per dividere una stringa in token separati da una o più caratteri (delimitatori).
\begin{lstlisting}[language=C]
	char* strtok(char* str, const char* delim);
\end{lstlisting}
Dove i parametri sono:
\begin{itemize}
	\item \textit{str} è la stringa da dividere
	\item \textit{delim} è la stringa di delimitatori, ogni carattere è un separatore tra i token
\end{itemize}
La funzione restituisce un puntatore al prossimo token o \textit{NULL} quando non ci sono più token. Inserisce $\backslash 0$ dopo ogni token nella stringa originale che quindi non è più utilizzabile.
\newpage
\noindent La versione thread-safe è \textbf{\textit{strtok\_r}}:
\begin{lstlisting}[language=C]
	strtok\_r(char* str, const char * delim, char ** saveptr);
\end{lstlisting}
Dove i parametri sono:
\begin{itemize}
	\item \textit{str} stringa da analizzare
	\item \textit{delim} stringa di delimitatori
	\item \textit{saveptr} puntatore che memorizza lo stato tra le chiamate
\end{itemize}
\subsection{calloc}
\textbf{Spiega la funzione \textit{calloc}.}\\
La funzione \textit{calloc} è utilizzata per allocare memoria dinamicamente. È simile a \textit{malloc} ma con inizializza la memoria allocata a $0$.
\begin{lstlisting}[language=C]
		void* calloc(size_t nitems, size_t size);
\end{lstlisting}
\begin{itemize}
	\item \textit{nitems} numero di elementi da allocare
	\item \textit{size} dimensione di ciascun elemento
\end{itemize}
Restituisce un puntatore alla memoria allocata.
\subsection{qsort}
\textbf{Come funziona \textit{qsort}? Quali sono le sue problematiche?}\\
\textit{qsort} è una funzione di libreria per ordinare un array di elementi.
\begin{lstlisting}[language=C]
	void qsort(void* base, size_t num, size_t size, int (*compar)(const void*, const void*));
\end{lstlisting}
Dove i parametri sono:
\begin{itemize}
	\item \textit{base} puntatore dell'array da ordinare
	\item \textit{num} numero di elementi nell'array
	\item \textit{size} dimensione di ciascun elemento nell'array
	\item \textit{compar} funzione di confronto personalizzata
\end{itemize}
La funzione utilizza \textit{void *}. Se il tipo non viene gestito correttamente possono verificarsi comportamenti imprevedibili. Se l'array contiene strutture la funzione \textit{compar} deve saperli gestire correttamente.
\subsection{fread e fwrite}
\textbf{Lettura e scrittura dei file binari con \textit{fread} e \textit{fwrite}.}\\
\begin{lstlisting}[language=C]
	size_t fread(void* ptr, size_t size, size_t count, FILE* stream);
	size_t fwrite(const void* ptr, size_t size, size_t count, FILE* stream);
\end{lstlisting}
Dove i parametri sono:
\begin{itemize}
	\item \textit{ptr} puntatore al buffer dove leggere o scrivere
	\item \textit{size} dimensione di ciascun elemento
	\item \textit{count} numero di elementi
	\item \textit{stream} puntatore al file aperto
\end{itemize}
Le funzioni restituiscono il numero di elementi letti o scritti. Se è inferiore a \textit{count} potrebbe significare la presenza di errori.
\subsection{fork}
\textbf{Creazione di processi con \textit{fork}.}\\
La chiamata alla funzione \textit{fork} genera un nuovo processo figlio duplicando il processo padre. Il padre deve sempre aspettare il figlio invocando \textit{wait}, altrimenti esso diventerà un processo zombie (non terminabile correttamente dal SO).
\begin{lstlisting}[language=C]
	pit_t fork(void);
\end{lstlisting}
\subsection{exit e wait}
\textbf{Terminazione di un processo figlio con \textit{exit} e \textit{wait}.}\\
\begin{itemize}
	\item \textit{exit}: la funzione termina l'esecuzione del processo chiamante, chiude tutti i file aperti, esegue la pulizia e restituisce un codice di uscita al SO
	\begin{lstlisting}[language=C]
		void exit(int status);
	\end{lstlisting}
	\item \textit{wait}: il processo genitore può utilizzare questa funzione (o in alternativa \textit{waitpid}) per attendere che uno o più processi figli terminino. Questo gli permette di recuperare il codice di uscita ed evitare che nascano processi zombie.
	\begin{lstlisting}[language=C]
		pid_t wait(int *status);
		pid_t wait_pid(pid_t pid, int* status, int options);
	\end{lstlisting}
	Dove i parametri sono:
	\begin{itemize}
		\item \textit{pid} PID del processo figlio da attendere
		\begin{itemize}
			\item $>0$ il PID specifico
			\item $-1$ qualsiasi figlio
			\item $0$ qualsiasi figlio nello stesso gruppo di processi
		\end{itemize}
		\item \textit{status} puntatore dove viene memorizzato il codice di uscita
		\item \textit{options} opzioni aggiuntive
	\end{itemize}
\end{itemize}
\subsection{execl}
\textbf{Spiega l'istruzione \textit{execl}.}\\
Serve ad eseguire un determinato file.
\begin{lstlisting}[language=C, mathescape]
	int execl(const char* path, const char* arg, $\ldots$, NULL);
\end{lstlisting}
Dove i parametri sono:
\begin{itemize}
	\item \textit{path} percorso del programma da eseguire
	\item \textit{arg} lista di argomenti passati al programma dove il primo è convenzionalmente il nome del programma
	\item \textit{NULL} la lista di argomenti deve terminare con NULL per segnalare la fine
\end{itemize}
\newpage
\subsection{shm\_open, ftruncate e mmap}
\textbf{Uso e significato di \textit{shm\_open}, \textit{ftruncate} e \textit{mmap}.}\\
\begin{itemize}
	\item \textit{shm\_open} serve a creare o aprire un oggetto di memoria condivisa
	\begin{lstlisting}[language=C]
		int shm_open(const char* nome, int oflag, mode_t mode);
	\end{lstlisting}
	Dove i parametri sono:
	\begin{itemize}
		\item \textit{nome} nome dell'oggetto nella memoria condivisa
		\item \textit{oflag} flag che indica se creare, leggere o scrivere (\textit{O\_CREAT}, \textit{O\_RDONLY} o \textit{O\_RDWR})
		\item \textit{mode} permessi di accesso
	\end{itemize}
	Ritorna un file descriptor associato all'oggetto di memoria condivisa.
	\item \textit{ftruncate} è usato per impostare la dimensione di un file o di un oggetto di memoria condivisa.
	\begin{lstlisting}[language=C]
		int ftruncate(int fd, off_t length);
	\end{lstlisting}
	Dove i parametri sono:
	\begin{itemize}
		\item \textit{fd} file descriptor ottenuto con \textit{shm\_open}
		\item \textit{length} nuova dimensione del file
	\end{itemize}
	Restituisce $0$ in caso di successo, $-1$ altrimenti.
	\item \textit{mmap} serve a mappare un file o una regione di memoria (e.g. oggetto creato da \textit{shm\_open}) nello spazio di indirizzamento del processo
	\begin{lstlisting}[language=C]
		void* nmap(void* addr, size_t length, int prot, int flags, int fd, off_t offset);
	\end{lstlisting}
	Dove i parametri sono:
	\begin{itemize}
		\item \textit{addr} indirizzo suggerito per il mapping, solitamente NULL
		\item \textit{length} dimensione della regione da mappare
		\item \textit{prot} permessi di accesso
		\item \textit{flags} specifica opzioni
		\item \textit{fd} file descriptor
		\item \textit{offset} offset inizio del file
	\end{itemize}
	Restituisce la memoria mappata o \textit{MAP\_FAILED} in caso di errore.
	\subsection{atexit}
	\textbf{Descrivi la funzione \textit{atexit}.}\\
	È utilizzata per registrare funzioni da eseguire automaticamente quando il programma termina (e.g. pulizia come chiusura di un file). 
	\begin{lstlisting}[language=C]
		int atexit(void(*func)(void));
	\end{lstlisting}
	\newpage
	\subsection{getopt}
	\textbf{Come funziona \textit{getopt}?}\\
	È una funzione per analizzare le opzioni e gli argomenti passati dalla riga di comando.
	\begin{lstlisting}[language=C]
		int getopt(int argc, const char* argv[], const char* opstring);
	\end{lstlisting}
	Dove i parametri sono:
	\begin{itemize}
		\item \textit{argc} numero di argomenti passati da riga di comando
		\item \textit{argv} array di stringhe contenenti i parametri passati
		\item \textit{opstring} una stringa che definisce le opzioni accettate. Ogni carattere rappresenta un'opzione, se è seguito da : vuol dire che richiede un argomento
	\end{itemize}
	Restituisce il carattere dell'opzione trovata, $-1$ quando tutte le opzioni sono state elaborate e $?$ se trova un'opzione non valida.
\end{itemize}
\subsection{pthread\_create e pthread\_join}
\textbf{Descrivi i prototipi di \textit{pthread\_create} e \textit{pthread\_join}.}\\
\begin{itemize}
	\item Creazione di thread
	\begin{lstlisting}[language=C]
		int pthread_create(pthread_t * thread, const pthread_attr_t* attr, void * (*start_routine)(void), void* arg);
	\end{lstlisting}
	Dove i parametri sono:
	\begin{itemize}
		\item \textit{thread} puntatore alla variabile che conterrà il thread creato
		\item \textit{attr} specifica gli attributi del thread (NULL di solito)
		\item \textit{(*start\_routine)(void)} funzione che rappresenta il punto di inizio del thread. Deve accettare void* come parametro e restituire void*
		\item \textit{arg} puntatore dell'argomento da passare alla funzione di inizio
	\end{itemize}
	Restituisce $0$ se il thread è stato creato con successo.
	\item Attesa thread
	\begin{lstlisting}[language=C]
		int pthread_join(pthread_t thread, void** retval);
	\end{lstlisting}
	Dove i parametri sono:
	\begin{itemize}
		\item \textit{thread} identificatore del thread
		\item \textit{retval} puntatore al valore restituito dalla funzione eseguita dal thread
	\end{itemize}
	Restituisce $0$ se il thread ha terminato con successo.
\end{itemize}
\subsection{signal}
\begin{itemize}
	\item \textit{signal} è utilizzato per la gestione di segnali
	\begin{lstlisting}[language=C]
		void* signal(int signum, void* handler(int));
	\end{lstlisting}
	Dove i parametri sono:
	\begin{itemize}
		\item \textit{signum} il codice del segnale da gestire
		\item \textit{handler} la funzione che gestisce il segnale. Deve prendere in input un intero che corrisponde al valore del segnale.
	\end{itemize}
\end{itemize}
\subsection{sleep}
\textbf{Attesa di segnali con \textit{sleep}.}\\
Una volta configurato il gestore di segnali, il programma può attendere utilizzando la funzione \textit{sleep}. Infatti, non appena il kernel riceve un segnale, interrompe l'attesa e chiama l'handler.
\subsection{kill}
\textbf{Invio di segnale dalla riga di comando con \textit{kill}.}\\
\begin{lstlisting}[language=BASH]
	kill -signal pid
\end{lstlisting}
Dove i parametri sono:
\begin{itemize}
	\item \textit{signal} nome o numero del segnale da inviare al processo
	\item \textit{pid} process ID a cui inviare il segnale
\end{itemize}

\begin{note}
	Utilizzare il parametro \textit{-l} per vedere l'elenco dei segnali e i loro numeri.
\end{note}

\subsection{sigwait}
\textbf{Attesa dei segnali con \textit{sigwait}.}\\
\textit{sigwait} è una funzione che sospende il thread chiamante fino a quando uno dei segnali specificati non viene ricevuto. Il segnale viene poi rimosso dalla coda dei segnali pendenti del processo e il controllo ritorna al thread chiamante, consentendo di gestire il segnale in modo sincrono.
\begin{lstlisting}[language=C]
	int sigwait(const sigset_t * set, int* sig);
\end{lstlisting}
Dove i parametri sono:
\begin{itemize}
	\item \textit{set} puntatore ad un insieme di segnali che il thread è disposto ad attendere
	\item \textit{sig} puntatore dove verrà memorizzato il segnale ricevuto
\end{itemize}
Restituisce $0$ se è corretto.
\subsection{sigqueue e sigwaitinfo}
\textbf{Invio di segnali con \textit{sigqueue} e \textit{sigwaitinfo}.}\\
\begin{itemize}
	\item \textit{sigqueue} consente di inviare segnali real time ad un thread specifico con dati personalizzati
	\begin{lstlisting}[language=C]
		int sigqueue(pthread_t thread, int sig, const union signal value);
	\end{lstlisting}
	Dove i parametri sono:
	\begin{itemize}
		\item \textit{thread} identificazione del thread destinatario
		\item \textit{sig} codice del segnale
		\item \textit{value} dato associato al segnale
	\end{itemize}
	\newpage
	\item \textit{sigwaitinfo} permette di attendere un segnale e ricevere le informazioni aggiuntive
	\begin{lstlisting}[language=C]
		int sigwaitinfo(const sigset_t * set, siginfo_t * info);
	\end{lstlisting}
	Dove i parametri sono:
	\begin{itemize}
		\item \textit{set} insieme di segnali da attendere
		\item \textit{info} struttura che riceve le informazioni del segnale. È  di tipo \textit{siginfo\_t} che contiene:
		\begin{itemize}
			\item \textit{si\_signo} numero del segnale
			\item \textit{si\_code} codice del segnale
			\item \textit{si\_value} dato personalizzato associato al segnale
		\end{itemize}
	\end{itemize}
\end{itemize}
\end{document}
